% Converted from reference/DB4OFlare (MadCap Flare, Java target)

\chapter{Product Philosophy}
\label{ref:product-philosophy}

db4o was originally developed by Versant Corporation, a developer of object database technology based in Redwood City, California. Versant has since discontinued db4o. It remains open source and is now maintained at \href{https://github.com/blakemcbride/db4o}{\nolinkurl{github.com/blakemcbride/db4o}}.

Versant's commercial object database technology, targeting extreme scale systems, is powering some of the world's most demanding applications for fortune 1000 companies in industries including:

\begin{itemize}
\item Telecommunications: Alcatel-Lucent, Deutsche Telecom, France Telecom, Ericsson, NEC, Nortel, Orange, Samsung, and more.
\item Finance: Financial Times, New York Stock Exchange, Dow Jones, Reuters, London Clearing House, Bank of America, and more,
\item Transportation: Sabre, GE Railways, BNS Railways, Galileo, and more.
\item Defense: Raytheon, Northrop Grumman, Lockheed, and more.
\item BioInformatics: Mayo Clinic, St. Jude medical, Eidogen, Science Factory, and more.
\end{itemize}

db4o has users and customers coming from 170 different countries, from Albania to Zimbabwe, and ranging from world class leaders like Boeing, Bosch, Intel, Ricoh and Seagate to a wide range of highly innovative start-up companies.

It is Versant's and the db4o team's mission to give developers a choice differentiated from relational approaches when it comes to object persistence and thus make their life a lot easier. There is no mapping! No mapping annotations or XML mapping meta data. The db4o database is designed to be a universal, affordable product platform that is easy to learn and use. Versant's open source dual-license business model combines the power of an open source development community with servicing commercial customers' needs for a predictable product roadmap, indemnification, single point of contact and full tech support with fast response times. For those requiring the super scale database capabilities, you can find the same easy to learn and implement solution in Versant's commercial products which have been in development for over a decade. This technology is also far more affordable that traditional relational database systems such as Oracle, Sybase, SQL Server, etc and to boot users also enjoy overall reductions in system footprints by as much as 50\% due to less indexing data, simpler design, zero mapping.

\section{Data Persistence}
\label{ref:product-philosophy-data-persistence}

Software programs using different data persistence technologies are an integral part of contemporary informational space. More than often such systems are implemented with the help of object-oriented programming language (Java, C\#, etc.) and a relational database management system (Oracle, MySQL, etc.). This implementation originally contains a mismatch between relational and object worlds, which is often called object/relational impedance mismatch. The essence of the problem is in the way the systems are designed. Object systems consist of objects and are characterized by identity, state, behavior, encapsulation. The relational model consists of tables, columns, rows and foreign keys and is described by relations, attributes and tuples.

The object-relational mismatch has become enormously significant with the total adoption of OO technology. This resulted in the rapid development of object-relational mappers (ORM), such as Hibernate, EclipseLink or Entity Framework. This solution "cures" the symptoms of the object relational mismatch by adding a layer into the software stack that automates the tedious task of linking objects to tables. However, this approach creates a huge drain on system performance, drives up software complexity and increases the burden on software maintenance, thus resulting in higher cost of ownership. While the mapper solution may be feasible in large, administered datacenter environments, it is prohibitive in distributed and zero-administration architectures such as those required for embedded databases in client software, mobile devices, middleware or real-time systems.

Significant side effects of the object relational mismatch manifest themselves in unnecessary system overhead with bloated footprints and runtime performance issues. Of course, there is also overall time to market delays due to poor developer productivity. The overhead still exists in ORM because under the covers, the runtime is still query driven. And, despite improvements in productivity for developers, incremental changes to your object models reek havoc during ORM schema evolution pitfalls. The more complicated your models are, the more problematic keeping changes in sync with the internal mapping.

Primary performance issues come from the fact that despite being called a "relational database", an RDBMS does not store direct relations. Relations are resolved at runtime by performing set based operations on primary-foreign key pairs. This means the application has to constantly re-discover data relationships at runtime resulting in immense CPU consumption for something that should be an inherent part of your application model. Further, because discovering these relations over and over again requires continual access to index structures and data to perform the set operations, contention is much higher within database internals leading to poor scalability of individual database processes.

Further, lack of direct storage of relations cause the application design to become query driven instead of object modeling driven. Using an object database, the relations are a fundamental part of the storage architecture. So, application design is model driven. You do not have to suffer any performance overhead for discovering an M-M relationship. The relationships are just there and immediately accessible to the requesting thread. This makes the internal structures much simpler and therefore less contention exists with data requests being isolated to data of interest instead of leveraging indexes or sequential scans. The result, individual processes become more scalable under concurrency.

Technology is ever changing and today there is a whole world of object database experts in the software community. Anyone who is an expert in ORM technology is an expert in object database technology. All of the concepts found in object life cycle management within ORM technologies were invented by the object database community in the early 90's. All of the tuning concepts of closure, fetch configurations, value vs reference types, light weight transactions - are concepts created by and applicable to object database technologies. Now with the growing popularity of object based design and the proliferation of ORM tools, thousands of developers are becoming experts in the object database API.

\section{OODBMS}
\label{ref:product-philosophy-oodbms}

The emergence of distributed data architectures - in networks, on clients and embedded in "smart" products such as cars or medical devices - is causing companies in an array of industries to look beyond traditional RDBMS technology and ORM for an improved way to deal with object persistence.

They are searching for a solution that can handle an enormous number of often complex objects, offer powerful replication and query capabilities, reduce development and maintenance costs and require minimum to zero administration overhead.

These requirements can be fulfilled by using an Object Oriented Database Management System (OODBMS). OODBMS provides an ideal match with object oriented environments like Java and .NET reducing the cost of development, support and versioning and hence overall system costs.

Using OODBMS in software projects also better supports modern Agile software engineering practices like:

\begin{itemize}
\item Continuous refactoring.
\item Agile modeling.
\item Continuous regression testing.
\item Configuration management.
\item Developer "sandboxes".
\end{itemize}

For more information about OODBMS technology refer to the \href{http://www.odbms.org/}{ODBMS.ORG website}.

\section{db4o Position}
\label{ref:product-philosophy-db4o-position}

The db4o database came to the market in 2004 with a goal to become the mainstream persistence architecture for embedded applications (in which the database is invisible to the end user) in general, and for mobile and embedded devices running on Java or .NET, in particular. Versant's vision for db4o is to become the affordable, dominant, open source persistence solution of object oriented developers of Java and .NET. In a very short time, the db4o team has achieved mainstream adoption with a fast growing user community currently boasting over 60,000 members. Community adoption is continually driven by db4o's efficient innovative technology, native queries, deployment in Java and .NET and its open source dual licensing business model.

The target environments for db4o are persistence architectures where there is no database administrator present and no RDBMS legacy, i.e. primarily on equipment, mobile and desktop clients, and in the middleware. Typical industries of db4o customers include transportation, communication, automation, medical sciences, industrial, consumer and financial applications, among many others.

Existing customers range from world-class leaders like Boeing, Bosch, Intel, Ricoh, and Seagate to a broad range of highly innovative start-up companies - in the Americas, EMEA, and Asia-Pacific.

As a client-side, embeddable database, db4o is particularly suited to be deployed in devices with embedded software.

\subsection*{Open Source}

db4o is open source, available under the \hyperref[ref:license]{GNU General Public License} (p.~\pageref{ref:license}). Versant originally offered it under a dual license model, as pioneered by MySQL: the GPL and the db4o Opensource Compatibility License (dOCL) for open source projects, and a commercial license for embedding db4o in for-profit products. The dOCL and the commercial license are no longer offered.

\subsection*{Success Drivers}

Open Source platform usage is one of the key factors of db4o success. db4o's openness attracted a vast (60,000 and counting) community of users and contributors. Through the community support db4o gets broad and immediate testing, receives constructive suggestions (from the users actually looking into the code) and invaluable peer exchange of experiences - positive and negative.

Another factor to db4o success is the technology used. As a new-generation object database, native to both Java and .NET, db4o eliminates the traditional trade-off between performance and object-orientation. Recent PolePosition benchmark results show that db4o outperforms object-relational mappers by orders of magnitude, up to 44x in use cases with complex object models.

db4o uniquely offers object persistence with zero-administration, object-oriented querying, replication and browsing capabilities, and a small footprint. Its single library (JAR/DLL) is easily deployed and runs in the same memory process as the application, making it a fully integrated and tunable portion of the developers application.

Customers, analysts, and experts agree that the db4o object database is one of the world's best and most popular choices, because it stores and retrieves objects natively and not only eliminates the overhead and resource consumption of an ORM, but also greatly reduces the product development and maintenance costs, resulting in a lean, fast and easily integratable into an OO development environment persistence solution, far superior in many cases to that of any RDBMS.

\subsection{db4o Applications}
\label{ref:product-philosophy-db4o-applications}

db4o can be used in a wide range of production and educational software. The primary focus is on embedded usage, like mobile systems (phones and handhelds), device electronics (printers, cars, robots), SCADA systems etc. The following table provides many (but not all) possible implementations with an explanations of the benefits of db4o in the selected environment:

{\small
\begin{longtable}{|p{\dimexpr(\linewidth-3\arrayrulewidth)*1/2-2\tabcolsep\relax}|p{\dimexpr(\linewidth-3\arrayrulewidth)*1/2-2\tabcolsep\relax}|}
\hline
\textbf{Environment} & \textbf{Benefits} \\ \hline
Educational systems & One-line persistence, Object-oriented model, intuitive programming interface make db4o an ideal educational tool. It is easy to use and it provides a meaningful example of object-oriented world. It is also native to most widely used OO languages: Java and .NET \\ \hline
Prototypes & Using db4o to build a prototype system is much quicker than using an RDBMS. In case of db4o you do not need to create a data model. Further there is no need to map your object model to the database. The general persistence mechanism is almost transparent and requires minimum effort to adapt to. Automatic refactoring allows rapid change of classes without the necessity to update the database. \\ \hline
SCADA & Using db4o in SCADA systems allows to achieve high performance in caching and replay of the events. Another benefit is a small footprint and easy integration with Java and .NET programming languages. db4o can also be run as a memory database, providing better performance though minimizing disk access. \\ \hline
Mobile applications & Mobile applications can benefit from in-process database, which requires zero-administration. Synchronization with the main server can be done with the help of dRS. Automatic refactoring can be another valuable factor, which allows to skip the job of updating the databases when a new version of object model is implemented. \\ \hline
Device applications & Device applications enjoy the same benefits as Mobile applications. In addition, smaller footprints can be achieved by using the minimal Micro edition. \\ \hline
Open-source software & GPL. Native to Java and .NET. Easily integrates with any Java and .NET open-source products. \\ \hline
Web-applications & open-source, reporting support from several Java open-source reporting frameworks and .NET reporting API \\ \hline
\end{longtable}}

However, other applications might not be well suited for db4o.

For example, in situations where you have increasing amounts of data ( over 10 Gigabytes ) and high concurrency ( over 20 concurrent users/processes ) along with your complex models. In these cases, a larger database system, such as the Versant object database, is likely a more appropriate choice. Versant's customer applications span a wide range of use including those exhibiting 1000's of current transactions ( 100's of thousands of concurrent tx per second ) to 100's of gigabytes with some Versant customers in the 25T+ sized database.

Another case is when you have simple and flat data model, primarily used for reporting. Simple table-like models of tuple records may be better supported by an RDBMS. In this case, adhoc data access would be more important to your application than well defined use cases using an object model. Typically this is complimented with the need let your users to be able to grab one of the plethora of commercial tools to poke at the database in an adhoc fashion.

\subsection{Scalability}
\label{ref:product-philosophy-scalability}

db4o is intended for embedded use with smaller databases around 2-16 GByte and at maximum 256 GByte. As a general rule, if you expect your database to grow beyond 16 gigabytes, you should look at a larger database system.

db4o is explicitly single-threaded. Concurrent accesses will be synchronized against a global database lock. That means db4o cannot deal with a highly concurrent access, since it blocks on all operations. When you expect a high load and concurrent access, then you should consider a larger database system.

If you want to make sure that your application can grow and scale with db4o database you may take the following steps:

\begin{itemize}
\item Write performance tests to accommodate for performance requirements of the growing application. Especially look that you capture the complexity of the model and don't use a simplified flat data model.
\item Also check that you can the throughput you need with concurrent access.
\end{itemize}

\section{Why Choose db4o}
\label{ref:product-philosophy-why-choose-db4o}

There are many advantages of using "native" object technology over RDBMS or RDBMS paired with an OR mapper, and these technological advantages significantly impact an organization's competitiveness and bottom line.

First, object databases not only simplify development by eliminating the resource-consuming OR-mismatch entirely, but they also foster more sophisticated and differentiated product development through gains in flexibility and productivity brought on by "true" object-orientation.

Second, with an object database, the object schema is the same as the data model. Developers can easier update their models to meet changing requirements, or for purposes of debugging or refactoring. db4o lets developers work with object structures almost as if they were "in-memory"' structures. Little additional coding is required to manage object persistence. As a result, companies can add new features to their products faster to stay ahead of the competition.

Third, developers can now use object-oriented and entirely native approaches when it comes to querying, since db4o was the first in the industry to provide Native Queries (NQ) with its Version 5.0 launched in November 2005 and since introduction of LINQ in .NET version 3.5. Db4o Native Queries provide an API that uses the programming languages Java or .NET to access the database. No time is spent on learning a separate data manipulation or query language. Unlike incumbent, string-based, non-native APIs (such as SQL, OQL, JDOQL and others) Native Queries and LINQ are 100\% type-safe, 100\% refactorable and 100\% object-oriented, boosting developer productivity by up to 30\%. In addition, the sophisticated modern programming development environments can be used to simplify the development and maintenance work even further.

Fourth, db4o's ground-breaking object-oriented replication technology solves problems arising from distributed data architectures. Partially connected devices need to efficiently replicate data with peers or servers. The challenge lies in the creation of "smart" conflict resolution algorithms, when redundant data sets are simultaneously modified and need to be merged. With db4o's OO approach, developers can build smarter and easier synchronization conflict resolution and embed the necessary business logic into the data layer, rather than into the middle-tier or application layer. This creates "smart" objects that can be stored in a distributed fashion, but easily consolidated, as the object itself knows how to resolve synchronization conflicts. It also enables db4o solution on a client to synchronize data with an RDBMS backend server.

As a result, developers can now more consistently persist data on distributed, partially connected clients than ever before, while decreasing bandwidth requirements and increasing the responsiveness and reach of their mobile solutions or smart devices to make products more competitive in the marketplace.

Fifth, db4o also allows for more complex object models than its relational or non-native counterparts do. As the persistence requirements become more complex, db4o's unique design easily handles (or absorbs) the added complexity, so developers can continue to work as though new complexity were never introduced. Complexity means not only taller object trees and extensive use of inheritance, but also dynamically evolving object models, most extremely if development is taking place under runtime conditions (which makes db4o a leading choice for biotech simulation software, for instance). db4o could be referred to as "agnostic to complexity," because it can automatically handle changes to the data model, without requiring extra work. No type or amount of complexity will change its behavior or restrain its capabilities, as is the case with RDBMS or non-native technology. With db4o breaking through this complexity, developers are able to write more user friendly and business-appropriate software components without incurring such high costs and modify them as needed, throughout the life cycle of the product with the same low cost.

In sum, db4o's native, object oriented architecture enables its users to build more competitive products with faster update cycles, more natural object models that match more realistically their use cases, and more distributed data architectures to increase the reach of products. db4o is clearly more flexible and powerful for embedded DB applications than any non-native OODBMS or RDBMS technologies available.

